\section{提交的结果文件与程序组织}

本文按题目要求提交四份结果文件，其行列规格与题面指定的输出间隔一致，全部数值保留四位小数。文件内容与规格列于表~\ref{tab:files}。

\begin{table}[H]
\centering
\caption{按题目要求提交的结果文件}\label{tab:files}
\begin{tabular}{llll}
\toprule
文件 & 对应问题 & 工作表与规格 & 时间与径向间隔 \\
\midrule
result1.xlsx & 问题一 & 温度、水分浓度两表，各 $1800\times21$ & $1$~s，$0.1$~cm \\
result2.xlsx & 问题二 & 温度、水分浓度两表，各 $10800\times21$ & $1$~s，$0.1$~cm \\
result3.xlsx & 问题三 & 水分浓度一表，$3452\times21$ & $60$~s，$0.1$~cm \\
result4.xlsx & 问题四 & 水分浓度一表，$3069\times22$ & $60$~s，$0.1$~cm \\
\bottomrule
\end{tabular}

\vspace{0.6em}
\begin{minipage}{0.94\textwidth}
\footnotesize 行数为数据行，不含表头；列数为径向位置数。四份文件依次覆盖 $1$ 至 $1800$~s、
$1$ 至 $10800$~s、$60$ 至 $207120$~s 与 $60$ 至 $184140$~s。第四份比其余多一列「药材表面」，
取 $s=1$ 物质节点的原值而不作空间插值；表内 $23191$ 格因该径向位置已落在移动边界之外而按题意留空，
不填零值。
\end{minipage}
\end{table}

程序按功能分为参数与工具、四问求解、校验与对比、结果导出四组，共十四个源文件，全部代码见附录~\ref{app:code}。参数模块只保存题面原值，任何派生量均在使用处现算，以避免同一参数出现两份不一致的副本；工具模块集中实现物性公式、界面系数的积分平均、控制体体积的精确积分与三对角求解；四个求解文件各自只负责本问特有的设置，公共推进逻辑不重复实现。
